\section{The Introductory Rites}

\subsection{What the opening is for}

The Missal gives the opening rites a single stated purpose, and it is worth having in front of
one before examining the parts: the rites preceding the Liturgy of the Word --- entrance,
greeting, penitential act, \latin{Kyrie}, \latin{Gloria}, and Collect --- have the character of
a beginning, an introduction, and a preparation, and their purpose is that the faithful coming
together should constitute a communion and dispose themselves rightly to hear God's word and
celebrate the Eucharist worthily.\endnote{\latin{IGMR} 46: ``Ritus qui liturgiam verbi
praecedunt, scilicet introitus, salutatio, actus paenitentialis, Kyrie, Gloria et collecta,
characterem habent exordii, introductionis et praeparationis. Finis eorum est, ut fideles in
unum convenientes communionem constituant et recte ad verbum Dei audiendum digneque
Eucharistiam celebrandam sese disponant.'' Read in the registered reproduction of the 2002
Latin.} That is a functional definition, and it explains the shape of what follows: a movement
inward (procession), a constitution of the body as addressed and answering (sign of the cross,
greeting), an acknowledgement of unworthiness, a cry to the Lord, praise, and a gathering of
all of it into one presidential petition.

\subsection{Entrance, altar, and the sign of the cross (nn.\ 1--2)}

The book's first rubric is about bodies in motion: with the people gathered, the priest with
the ministers goes to the altar while the entrance chant is carried out.\endnote{\latin{Ordo
Missae} 1: ``Populo congregato, sacerdos cum ministris ad altare accedit, dum cantus ad
introitum peragitur.'' The same number directs the profound bow, the kiss of the altar, the
optional incensation of cross and altar, and the going to the chair.} The order of clauses
matters. The people are already gathered when the rite begins; the ministers come to them and
pass through them; and the destination of the procession is not the chair but the altar, which
is reverenced with a bow and kissed before the priest goes to preside. The altar is thus
greeted before the assembly is, and the first act of the celebrant's body is an act of homage
to a stone. Whatever else is true of this Order, its opening gesture is not anthropocentric.

The entrance chant has a stated function --- to open the celebration, foster the unity of those
gathered, introduce their minds to the mystery of the liturgical time or festivity, and
accompany the procession --- and in the United States it may be drawn from four sources: the
antiphon in the Missal, or the antiphon with its psalm from the \latin{Graduale Romanum}; the
antiphon with psalm from the \latin{Graduale Simplex}; a chant from another approved collection
of psalms and antiphons; or another suitable liturgical chant approved by the Conference or the
diocesan Bishop.\endnote{\latin{IGMR} 47--48; the four options are the United States adaptation
of n.\ 48, read at the official USCCB presentation of Chapter II, 2026-07-25. The parallel
adaptation for the Communion chant is at n.\ 87 and is discussed in its own place. The Latin
typical text of n.\ 48 does not contain the four-option list, which is territorial.} That
fourth option is the widest door in the Introductory Rites and the one through which most of
what an American Catholic actually hears at the entrance has come; it is a territorial
adaptation, not a provision of the typical edition, and it is worth naming as such rather than
attributing the resulting repertoire to the Missal.

The sign of the cross is made by priest and people together, standing, while the priest, turned
toward the people, says the trinitarian formula and the people answer \latin{Amen}. This is a
postconciliar arrangement in its present shape: in the older Order the priest signed himself at
the foot of the altar during the prayers preparatory to Mass, and the assembly's corporate
signing at the outset is a restoration of visibility to what had become the celebrant's private
beginning.

\subsection{The greeting and what the answer concedes (n.\ 2)}

Three forms are printed --- the Pauline trinitarian greeting from 2~Corinthians 13:13, the
shorter Pauline greeting, and \latin{Dominus vobiscum}; a bishop says \latin{Pax vobis} in this
first greeting instead.\endnote{\latin{Ordo Missae} 2, giving \latin{Gratia Domini nostri Iesu
Christi, et caritas Dei, et communicatio Sancti Spiritus sit cum omnibus vobis}; \latin{Gratia
vobis et pax a Deo Patre nostro et Domino Iesu Christo}; \latin{Dominus vobiscum}; and, for a
bishop, \latin{Pax vobis}.} The people's answer to all three is the same: \latin{Et cum spiritu
tuo}.

The English of that answer changed in 2011, and the reason is documented rather than
conjectural. \emph{Liturgiam authenticam} names this very response as an example of expressions
belonging to the heritage of the whole or a great part of the ancient Church which are to be
respected by a translation that is as literal as possible.\endnote{\emph{Liturgiam
authenticam} 56, read 2026-07-25: ``Certain expressions that belong to the heritage of the
whole or of a great part of the ancient Church, as well as others that have become part of the
general human patrimony, are to be respected by a translation that is as literal as possible,
as for example the words of the people's response \latin{Et cum spiritu tuo}, or the expression
\latin{mea culpa, mea culpa, mea maxima culpa} in the Act of Penance of the Order of Mass.''
Both examples are named in the instruction's own text; both were in fact changed in the 2011
English.} The 1973 English had rendered the answer as a reciprocal greeting; the 2011 English
gives ``and with your spirit,'' restoring both the noun and the ambiguity of the Latin. The
ambiguity is the point of the debate. Patristic and modern readings differ over whether
\latin{spiritus tuus} means simply ``you,'' after a Semitic idiom, or refers to the grace of
office conferred by ordination --- the second reading being why the response is reserved to
exchanges with an ordained minister. This study does not settle that question; it records that
the Latin admits both, that the older English closed it, and that the instruction's stated
reason for reopening it was neither theory: it was the antiquity and universality of the
formula.\endnote{The claim about the reservation of the dialogue to ordained ministers is a
matter of the Missal's own distribution of the formula, not of a doctrinal note in the
Instruction. This study makes no claim about which patristic interpretation is correct, and it
did not examine the patristic dossier on \latin{spiritus tuus}.}

\subsection{The Penitential Act (nn.\ 4--6)}

Three forms are printed, all introduced by the same invitation, \latin{Fratres, agnoscamus
peccata nostra, ut apti simus ad sacra mysteria celebranda}, and all followed by a brief pause
of silence. Form I is the \latin{Confiteor}, said by all together, with the striking of the
breast at \latin{mea culpa, mea culpa, mea maxima culpa}. Form II is a versicle-and-response
pair drawn from the psalms (\latin{Miserere nostri, Domine} --- \latin{Quia peccavimus tibi};
\latin{Ostende nobis, Domine, misericordiam tuam} --- \latin{Et salutare tuum da nobis}). Form
III sets tropes before the \latin{Kyrie} acclamations, of which the Missal prints one set and
permits others; the tropes are addressed to Christ throughout, which is why in this form the
\latin{Kyrie} that follows is not repeated. All three end with the priest's absolution formula,
which the General Instruction says lacks the efficacy of the sacrament of
Penance.\endnote{\latin{Ordo Missae} 4--7; \latin{IGMR} 51, which says of the priest's
concluding formula that it ``efficacia sacramenti Paenitentiae caret.'' Note the small but
consequential rubric at n.\ 7: the \latin{Kyrie} follows ``nisi iam praecesserint in aliqua
formula actus paenitentialis'' --- that is, unless it has already been used in form III.} A
footnote on the same page provides that on Sundays, especially in Eastertide, the blessing and
sprinkling of water in memory of baptism may replace the customary penitential act.

Two points of history are secure and worth separating. The \latin{Confiteor} itself is medieval
in its developed form and was, in the older Roman Order, part of the prayers at the foot of the
altar --- said by the priest and answered by the ministers, that is, a clerical dialogue.
Making it the assembly's act is a change of speaker, not an invention of a text. The
penitential act as a distinct opening moment, by contrast, is presented by Paul VI's
constitution as a restoration: among the things restored ``to the ancient norm of the holy
Fathers'' the constitution names the homily, the universal prayer, and ``the penitential rite,
or act of reconciliation with God and with the brethren, to be performed at the beginning of
Mass,'' to which its due importance has been restored.\endnote{Paul VI, apostolic constitution
\latin{Missale Romanum} (3 April 1969), Latin as reprinted in the third typical edition and
read publication-locally 2026-07-25: ``Huc accedit quod restituuntur \ldots\ ad pristinam
sanctorum Patrum normam nonnulla quae temporum iniuria deciderunt; cuius generis sunt homilia
et oratio universalis seu oratio fidelium et ritus paenitentialis, seu reconciliationis cum Deo
et cum fratribus, initio Missae peragendus: cui, ut oportebat, suum redditum est momentum.''}
The word \latin{restituuntur} is a historical claim made by the promulgating authority, and a
reader is entitled to weigh it: what the reform restored here was a penitential moment
belonging to the whole assembly at the threshold of the rite, which is genuinely ancient in
type; the particular texts used to fill it are not all ancient, and form II and form III are
new compositions in their present arrangement.

The 2011 English of the \latin{Confiteor} restored the triple \latin{mea culpa}, which the 1973
English had reduced to a single admission. \emph{Liturgiam authenticam} had named that very
phrase, alongside \latin{Et cum spiritu tuo}, as an expression to be translated as literally as
possible. The English text itself is not reproduced here; what can be said exactly is that the
triple form is in the Latin, that the instruction required it, and that the 2011 text has it.

\subsection{\latin{Kyrie} and \latin{Gloria} (nn.\ 7--8)}

The \latin{Kyrie} is printed in Greek and in the sixfold form, two acclamations to each
invocation, with the note that other melodies are found in the \latin{Graduale
Romanum}.\endnote{\latin{Ordo Missae} 7. The Missal prints \latin{V. Kyrie, eleison. R. Kyrie,
eleison.} and so on --- three invocations, each doubled --- rather than the ninefold form of
the older Roman use.} The retention of the Greek in a Latin book, and of a Latin-and-Greek
acclamation in a vernacular translation, is the clearest single instance of the principle that
some formulas are kept because of what their survival means rather than because of what they
say. The sixfold rather than ninefold arrangement is a simplification of the older Roman
practice and is presented in the General Instruction as an acclamation by which the faithful
acclaim the Lord and implore his mercy, ordinarily done by all with a part for the choir or a
cantor.

The \latin{Gloria} is printed complete and is described as a very ancient and venerable hymn by
which the Church, gathered in the Holy Spirit, glorifies and entreats God the Father and the
Lamb.\endnote{\latin{Ordo Missae} 8 prints the hymn; \latin{IGMR} 53 describes it and rules its
use: sung or said on Sundays outside Advent and Lent, on solemnities and feasts, and at
particular more solemn celebrations. The Latin text at n.\ 8 is unchanged from the older Roman
form.} Its conditional status is worth naming precisely, because the calendar does most of the
work: its absence in Advent and Lent is not a penitential improvisation but the reason the
Sundays of those seasons sound different from the Sundays that surround them.

\subsection{The Collect (n.\ 9)}

The Collect is the first of the three presidential orations and the structural close of the
opening. The rubric is unusually exact about what happens before the prayer: with hands joined
the priest says \latin{Oremus}, and all, together with the priest, pray in silence for some
space of time; then, with hands extended, he says the Collect, at whose end the people acclaim
\latin{Amen}.\endnote{\latin{Ordo Missae} 9: ``Quo hymno expleto, sacerdos, manibus iunctis,
dicit: \latin{Oremus}. Et omnes una cum sacerdote per aliquod temporis spatium in silentio
orant. Tunc sacerdos, manibus extensis, dicit orationem collectam, qua expleta, populus
acclamat: \latin{Amen}.''} The silence is prescribed, not permitted; \latin{Oremus} is an
invitation to the assembly to pray, and the oration that follows collects what has been prayed.
This is the plainest sense of the traditional name, and the rubric preserves the mechanism that
gives the name its meaning. Where the silence is omitted --- a common shortcut --- the priest's
oration ceases to collect anything, and the movement's own logic is lost.

The oration itself is proper to the day and therefore outside the scope of this study, which
treats the Order of Mass rather than the temporal and sanctoral cycles. Two of its formal
features belong here. It is addressed, as the great majority of Roman orations are, to the
Father; and it concludes through Christ, in the unity of the Holy Spirit, so that the
assembly's petition is offered through the one Mediator --- a structure the General Instruction
underlines when it says that the priest addresses the prayer to God in the name of the whole
holy people and of all present, and that the people, joining themselves to the petition, make
the prayer their own by the acclamation \latin{Amen}.

The 2011 English of the orations is a study in itself and cannot be conducted here with
quotation, because the approved English of the Collects is under copyright and this study does
not reproduce it. What can be said exactly, and checked by any reader with both books, is
structural: the 1973 English generally broke the long Latin period into short independent
sentences and reduced the vocative and relative clauses; the 2011 English generally preserves
the period, the subordination, and the distinct address, in obedience to an instruction that
expressly required the Roman Rite's concise and compact manner of expression to be maintained
insofar as possible, and that varied Latin vocabulary should give rise to corresponding variety
in the vernacular.\endnote{\emph{Liturgiam authenticam} 51 and 57, read 2026-07-25. N.\ 51
names, as examples of impoverishing uniformity, the rendering of \latin{satiari, sumere,
vegetari, pasci} by one word, and of \latin{Domine, Deus, Omnipotens aeterne Deus, Pater} by one
form of address. The characterization of the two English translations' handling of the Latin
period is this study's own structural judgement, offered as such, and is not a quotation of
either text.}
